\begin{lemma}[Parity of multidimensional Cayley--Poisson modes]%
\label{lem:multidim-cayley-poisson-mode-parity}
For every \(d\ge1\) and every multiindex \(\alpha\in\NN_0^d\),
\begin{equation}\label{eq:multidim-cayley-poisson-mode-parity}
U_\alpha^{(d)}(-y)=(-1)^{|\alpha|}U_\alpha^{(d)}(y)
\qquad(y\in\RR^d).
\end{equation}
Consequently, for any multiindices
\(\alpha_1,\ldots,\alpha_r\),
\begin{equation}\label{eq:multidim-cayley-poisson-product-parity}
\int_{\RR^d}\prod_{j=1}^r
U_{\alpha_j}^{(d)}(y)\,\Omega_d(\dd y)=0
\qquad\text{if }\ |\alpha_1|+\cdots+|\alpha_r|\text{ is odd}.
\end{equation}
\end{lemma}

\begin{proof}
The defining quotient satisfies
\[
F_{-y}(z)
=\left(\frac{1+|y|^2}{1+|-y-z|^2}\right)^{(d+1)/2}
=F_y(-z).
\]
Both sides are real analytic in \(z\) near \(0\).  Comparing the Taylor
coefficients in
\[
F_{-y}(z)=\sum_{\alpha}U_\alpha^{(d)}(-y)z^\alpha,
\qquad
F_y(-z)=\sum_{\alpha}(-1)^{|\alpha|}U_\alpha^{(d)}(y)z^\alpha
\]
gives \eqref{eq:multidim-cayley-poisson-mode-parity}.  Multiplying this
identity for \(\alpha_1,\ldots,\alpha_r\) gives
\[
\prod_{j=1}^r U_{\alpha_j}^{(d)}(-y)
=(-1)^{|\alpha_1|+\cdots+|\alpha_r|}
\prod_{j=1}^r U_{\alpha_j}^{(d)}(y).
\]
Since \(\Omega_d(\dd y)=c_d(1+|y|^2)^{-(d+1)/2}\dd y\) is invariant under
\(y\mapsto-y\), the integral of the product is equal to its negative when
the total order is odd, proving
\eqref{eq:multidim-cayley-poisson-product-parity}.
\end{proof}

\medskip
\noindent\textbf{Deterministic quotient remainder.}
The next lemma is the analytic estimate used later to verify the named
matched quotient jet from finite moments.  The quotient jet itself, including
the entropy-existence clause, is defined in full immediately before the
transfer proposition so the analytic input is visible at the point where it is
used.

\begin{lemma}[Global multidimensional quotient remainder]%
\label{lem:global-multidim-quotient-remainder-core}
Let \(d\ge1\), \(a=(d+1)/2\), and
\[
F_y(z)=\left(\frac{1+|y|^2}{1+|y-z|^2}\right)^a .
\]
For a multiindex \(\alpha\), put
\[
U_\alpha^{(d)}(y)=\frac1{\alpha!}\partial_z^\alpha F_y(0).
\]
If \(L\in\mathbb N\) and \(L\ge d+1\), then there is a constant
\(C_{d,L}<\infty\), depending only on \(d\) and \(L\), such that
\begin{equation}\label{eq:multidim-coefficient-bound}
\sup_{y\in\RR^d}|U_\alpha^{(d)}(y)|\le C_{d,L}
\qquad (|\alpha|\le L),
\end{equation}
\begin{equation}\label{eq:multidim-small-translation-remainder}
\sup_{y\in\RR^d}\left|
F_y(z)-\sum_{|\alpha|\le L}U_\alpha^{(d)}(y)z^\alpha
\right|
\le C_{d,L}|z|^{L+1}
\qquad (|z|\le1),
\end{equation}
and
\begin{equation}\label{eq:multidim-large-translation-remainder}
\sup_{y\in\RR^d}\left|
F_y(z)-\sum_{|\alpha|\le L}U_\alpha^{(d)}(y)z^\alpha
\right|
\le C_{d,L}|z|^{L}
\qquad (|z|>1).
\end{equation}
Equivalently, for all \(y,z\),
\begin{equation}\label{eq:multidim-global-translation-remainder}
\left|
F_y(z)-\sum_{|\alpha|\le L}U_\alpha^{(d)}(y)z^\alpha
\right|
\le C_{d,L}\left(
|z|^{L+1}{\bf1}_{\{|z|\le1\}}
{}
+{}|z|^L{\bf1}_{\{|z|>1\}}
\right).
\end{equation}
Consequently, if \(X\) is an \(\RR^d\)-valued random vector with
\(\EE|X|^L<\infty\), \(\Omega_d\) is the Cauchy probability measure above,
and
\[
r_t(y):=\EE\left[
F_y(X/t)-\sum_{|\alpha|\le L}U_\alpha^{(d)}(y)(X/t)^\alpha
\right],
\]
then
\begin{equation}\label{eq:multidim-integrated-remainder-bound}
\|r_t\|_{L^\infty(\RR^d)}+\|r_t\|_{L^1(\Omega_d)}
\le C_{d,L}t^{-L}\EE\left[
|X|^L\frac{|X|}{t}{\bf1}_{\{|X|\le t\}}
{}
+{}|X|^L{\bf1}_{\{|X|>t\}}
\right]=o(t^{-L}).
\end{equation}
\end{lemma}

\begin{proof}
The coefficient bound \eqref{eq:multidim-coefficient-bound} and the
small-translation remainder
\eqref{eq:multidim-small-translation-remainder} are exactly
Lemma~\ref{lem:local-multidim-quotient-taylor-remainder-core}.  Thus the
constant in these two estimates depends only on \(d,L\).

It remains to handle large translations.  Fix \(|z|>1\).  First suppose
\(|y-z|\le |z|/2\).  Then \(|y|\le3|z|/2\), and hence
\[
F_y(z)\le (1+|y|^2)^a\le C_d|z|^{d+1}.
\]
On the complementary region \(|y-z|>|z|/2\), the elementary estimate
\[
1+|y|^2
\le 2(1+|y-z|^2+|z|^2)
\le 2(1+|z|^2)(1+|y-z|^2)
\]
gives
\[
F_y(z)
\le 2^a(1+|z|^2)^a
\le C_d|z|^{d+1}
\qquad(|z|>1).
\]
Thus, uniformly in \(y\),
\[
\sup_y F_y(z)\le C_d|z|^{d+1}\le C_{d,L}|z|^L
\qquad(|z|>1),
\]
where the last inequality is precisely where \(L\ge d+1\) is used.  The
coefficient bound from the local lemma gives, again uniformly in \(y\),
\[
\sum_{|\alpha|\le L}|U_\alpha^{(d)}(y)|\,|z|^{|\alpha|}
\le C_{d,L}\sum_{k=0}^{L}|z|^k
\le C_{d,L}'|z|^L
\qquad(|z|>1).
\]
Consequently
\[
\sup_y\left(
|F_y(z)|
+\sum_{|\alpha|\le L}|U_\alpha^{(d)}(y)|\,|z|^{|\alpha|}
\right)
\le C_{d,L}|z|^L
\qquad(|z|>1),
\]
which proves \eqref{eq:multidim-large-translation-remainder}.  Combining
this estimate with the small-translation remainder proves the global form
\eqref{eq:multidim-global-translation-remainder}.

Set \(z=X/t\) in \eqref{eq:multidim-global-translation-remainder}.  Taking
expectations gives the pointwise bound, uniform in \(y\),
\[
|r_t(y)|
\le
C_{d,L}\left[
t^{-(L+1)}
\EE\!\left(|X|^{L+1}{\bf1}_{\{|X|\le t\}}\right)
+t^{-L}
\EE\!\left(|X|^L{\bf1}_{\{|X|>t\}}\right)
\right].
\]
Equivalently,
\[
t^{-(L+1)}
\EE\!\left(|X|^{L+1}{\bf1}_{\{|X|\le t\}}\right)
=
t^{-L}
\EE\!\left[
|X|^L\frac{|X|}{t}{\bf1}_{\{|X|\le t\}}
\right]
=o(t^{-L}).
\]
because \((|X|/t){\bf1}_{\{|X|\le t\}}\le1\) and converges pointwise to
zero, while \(|X|^L\) is integrable.  Also
\[
t^{-L}\EE\!\left(|X|^L{\bf1}_{\{|X|>t\}}\right)=o(t^{-L})
\]
by the absolute continuity of the integral for the integrable function
\(|X|^L\).  These two estimates are independent of \(y\), and therefore give
the \(L^\infty\) part of \eqref{eq:multidim-integrated-remainder-bound}.
The \(L^1(\Omega_d)\) part follows from the same uniform bound because
\(\Omega_d\) is a probability measure.
\end{proof}

\medskip
\noindent\textbf{Matched multidimensional quotient input.}
For \(d\ge1\) and \(L\in\mathbb N\) with \(L\ge2\), put
\[
h_t=P_t^{(d)}*\nu,\qquad
g_t=P_t^{(d)}(\cdot-\bar{\gamma}),\qquad
\delta_t^{(d)}(y)=
\frac{h_t(\bar{\gamma}+ty)}{g_t(\bar{\gamma}+ty)}-1,
\]
and \(\Phi(s)=(1+s)\log(1+s)-s\).

\begin{definition}[Large-\(t\) entropy-existence clause]%
\label{def:multidim-ex-dl}
For the above \(h_t,g_t,\delta_t^{(d)}\), the condition
\((\mathrm{EX}_{d,L})\) means that there is \(T<\infty\) such that, for every
\(t\ge T\),
\[
1+\delta_t^{(d)}(y)>0,\qquad
\int_{\RR^d}\delta_t^{(d)}(y)\,\Omega_d(\dd y)=0,\qquad
\int_{\RR^d}\Phi(\delta_t^{(d)}(y))\,\Omega_d(\dd y)<\infty,
\]
and that the matched KL integral is represented by
\[
D_{\rm KL}(h_t\|g_t)
=\int_{\RR^d}\Phi(\delta_t^{(d)}(y))\,\Omega_d(\dd y).
\]
The symbol records only these large-\(t\) positivity, zero-mass,
integrability, and change-of-variables assertions; it is not a moment
hypothesis and it is not implied by the finite-covariance algebra below.
\end{definition}

We say that \((\nu,d,L)\) satisfies the matched Poisson quotient jet
\((\mathrm{MPJ}_{d,L})\) when \(\nu\) is a probability law on \(\RR^d\), the
mean \(\bar{\gamma}\) exists, the centred tensor moments
\[
\mu_\alpha=\int_{\RR^d}(\gamma-\bar{\gamma})^\alpha\,\dd\nu(\gamma)
\qquad(2\le|\alpha|\le L)
\]
exist, and the scaled quotient has the \(L^\infty\cap L^1(\Omega_d)\) jet
\begin{equation}\label{eq:mpj-dl-definition}
\begin{aligned}
(\mathrm{MPJ}_{d,L}):\qquad
\delta_t^{(d)}(y)
&=
\sum_{2\le\lvert\alpha\rvert\le L}
\mu_\alpha U_\alpha^{(d)}(y)t^{-\lvert\alpha\rvert}
+\rho_{L,t}^{(d)}(y),\\
\lVert\rho_{L,t}^{(d)}\rVert_\infty+
\lVert\rho_{L,t}^{(d)}\rVert_{L^1(\Omega_d)}
&=o(t^{-L}),\qquad
\lVert\delta_t^{(d)}\rVert_\infty=O(t^{-2}).
\end{aligned}
\end{equation}
Thus the input used in
Proposition~\ref{thm:multidimensional-poisson-entropy-jets} has two parts:
the quotient jet \((\mathrm{MPJ}_{d,L})\) and the entropy-existence clause
\((\mathrm{EX}_{d,L})\).  The restriction \(L\ge d+1\) enters when the
large-translation estimates are closed by the finite \(L\)-moment remainder
estimate below.

\subsubsection*{Multidimensional proof route}\leavevmode\par\smallskip
The multidimensional proof has one analytic bridge.  For \(L\ge d+1\),
\[
\EE|X_c|^L<\infty
\quad\Longleftrightarrow\quad
\mathcal R_{d,L}(t)=o(t^{-L})
\quad\Longrightarrow\quad
(\mathrm{MPJ}_{d,L})+(\mathrm{EX}_{d,L}),
\]
and then Proposition~\ref{thm:multidimensional-poisson-entropy-jets} gives the
KL tensor jet.  At order four this jet has
\[
\mathcal A_4^{(d)}(\nu)
=\frac12\int_{\RR^d}(D_2^{(d)})^2\,\dd\Omega_d
=\mathcal Q_d(\Sigma),
\]
where the last equality is the covariance algebra proved below.  Thus finite
\(L\)-moment supplies the quotient/existence bridge and hence the actual
leading KL coefficient.  Finite covariance by itself supplies only the
algebraic quantity \(\mathcal Q_d(\Sigma)\).

The equivalence displayed in this proof route refers only to the functional
\(\mathcal R_{d,L}\) defined next.  It is not a statement that every
two-region estimate, or every possible quotient remainder norm, is equivalent
to finite \(L\)-moment.

\begin{definition}[Two-region remainder functional]%
\label{def:two-region-remainder-functional}
Let \(d\ge1\), let \(L\in\mathbb N\) satisfy \(L\ge d+1\), and let \(\nu\) be
a probability measure on \(\RR^d\) with mean \(\bar{\gamma}\).  Put
\(X_c=\gamma-\bar{\gamma}\).  For \(t>0\), define \(\mathcal R_{d,L}(t)\) by
\begin{equation}\label{eq:conditional-multidim-two-region-tail}
\begin{aligned}
\mathcal R_{d,L}(t)
&:=t^{-(L+1)}\EE\!\left[
|X_c|^{L+1}{\bf1}_{\{|X_c|\le t\}}\right] \\
&\quad{}+ t^{-L}\EE\!\left[
|X_c|^L{\bf1}_{\{|X_c|>t\}}\right],
\end{aligned}
\end{equation}
\end{definition}

\begin{lemma}[Two-region remainder is finite \(L\)-moment]%
\label{lem:two-region-remainder-equivalent-finite-moment}
Let \(d\ge1\), let \(L\in\mathbb N\) satisfy \(L\ge d+1\), and let \(\nu\) be
a probability measure on \(\RR^d\) with mean \(\bar{\gamma}\).  Put
\(X_c=\gamma-\bar{\gamma}\).  Then
\[
\mathcal R_{d,L}(t)=o(t^{-L})
\qquad(t\to\infty)
\]
is equivalent to \(\EE|X_c|^L<\infty\), with the convention that
\(\mathcal R_{d,L}(t)<\infty\) for all sufficiently large \(t\).
This is the only two-region equivalence used in the multidimensional
finite-moment theorem.
\end{lemma}

\begin{proof}
Assume first that \(\mathcal R_{d,L}(t)=o(t^{-L})\).  Choose \(t_0\) so large
that the terms in \eqref{eq:conditional-multidim-two-region-tail} are finite.
The bounded part \(\EE[|X_c|^L{\bf1}_{\{|X_c|\le t_0\}}]\) is finite, and
the second term gives
\[
\EE[|X_c|^L{\bf1}_{\{|X_c|>t_0\}}]<\infty .
\]
Thus \(\EE|X_c|^L<\infty\).

Conversely assume \(\EE|X_c|^L<\infty\).  The local term satisfies
\[
t^{-(L+1)}\EE[|X_c|^{L+1}{\bf1}_{\{|X_c|\le t\}}]
=t^{-L}\EE\!\left[
|X_c|^L\frac{|X_c|}{t}{\bf1}_{\{|X_c|\le t\}}
\right]
=o(t^{-L})
\]
by dominated convergence.  The tail term is
\[
t^{-L}\EE[|X_c|^L{\bf1}_{\{|X_c|>t\}}]=o(t^{-L})
\]
by tail integrability.
\end{proof}

The restriction \(L\ge d+1\) enters the deterministic large-translation
estimates through \(\mathcal R_{d,L}\).  The next proposition is the explicit
remainder estimate used by the finite-moment theorem; by
Lemma~\ref{lem:two-region-remainder-equivalent-finite-moment} it introduces
no hypothesis beyond \(\EE|X_c|^L<\infty\) in the present range.

\begin{proposition}[Explicit two-region quotient remainder estimate]%
\label{thm:multidimensional-conditional-tail-density-quotient}
Let \(d\ge1\), let \(L\in\mathbb N\) satisfy \(L\ge d+1\), and let \(\nu\) be
a probability measure on \(\RR^d\) with mean \(\bar{\gamma}\).  Put
\(X_c=\gamma-\bar{\gamma}\).  Assume
\[
\mathcal R_{d,L}(t)=o(t^{-L}).
\]
Equivalently, by Lemma~\ref{lem:two-region-remainder-equivalent-finite-moment},
this is exactly the finite \(L\)-moment condition
\(\EE|X_c|^L<\infty\).  Define
\[
\begin{gathered}
\mu_\alpha=\int_{\RR^d}X_c^\alpha\,\dd\nu(\gamma)
\qquad(|\alpha|\le L).
\end{gathered}
\]
Then, for
\[
\delta_t^{(d)}(y):=
\frac{(P_t^{(d)}*\nu)(\bar{\gamma}+ty)}
{P_t^{(d)}(ty)}-1,
\]
\begin{equation}\label{eq:multidim-density-quotient-expansion-main}
\delta_t^{(d)}(y)=
\sum_{2\le|\alpha|\le L}\mu_\alpha U_\alpha^{(d)}(y)t^{-|\alpha|}
+\rho_{L,t}^{(d)}(y),
\end{equation}
where the remainder satisfies the explicit \(L^\infty\cap L^1(\Omega_d)\)
bound
\begin{equation}\label{eq:multidim-density-quotient-explicit-remainder}
\begin{aligned}
\|\rho_{L,t}^{(d)}\|_\infty+
\|\rho_{L,t}^{(d)}\|_{L^1(\Omega_d)}
&\le C_{d,L}\mathcal R_{d,L}(t)  \\
&=o(t^{-L}).
\end{aligned}
\end{equation}
Moreover \(\|\delta_t^{(d)}\|_\infty=O(t^{-2})\), so
\((\mathrm{EX}_{d,L})\) holds for all sufficiently large \(t\).
\end{proposition}

\begin{proof}
Put \(h_t=P_t^{(d)}*\nu\) and \(g_t=P_t^{(d)}(\cdot-\bar{\gamma})\).
For \(x=\bar{\gamma}+ty\),
\[
\frac{P_t^{(d)}(x-\gamma)}{P_t^{(d)}(x-\bar{\gamma})}
=
\left(\frac{t^2+|x-\bar{\gamma}|^2}{t^2+|x-\gamma|^2}\right)^{(d+1)/2}
=F_y(X_c/t).
\]
Thus
\[
\delta_t^{(d)}(y)=\int_{\RR^d}\{F_y(X_c/t)-1\}\,\dd\nu(\gamma).
\]
Subtract the Taylor polynomial through order \(L\).  The remainder
\[
\rho_{L,t}^{(d)}(y):=
\int_{\RR^d}\left\{
F_y(X_c/t)-\sum_{|\alpha|\le L}U_\alpha^{(d)}(y)(X_c/t)^\alpha
\right\}\dd\nu(\gamma)
\]
satisfies \eqref{eq:multidim-density-quotient-explicit-remainder}.  Indeed,
on \(|X_c|\le t\), Lemma~\ref{lem:global-multidim-quotient-remainder-core}
gives the uniform Taylor remainder
\[
O\!\left(\left|\frac{X_c}{t}\right|^{L+1}\right).
\]
On \(|X_c|>t\), the large-translation part of the proof of that lemma gives
\(\sup_yF_y(X_c/t)\le C_dt^{-(d+1)}|X_c|^{d+1}\), while the coefficient bound
\eqref{eq:multidim-coefficient-bound} gives
\[
\sup_y\sum_{|\alpha|\le L}
|U_\alpha^{(d)}(y)|\,|X_c/t|^{|\alpha|}
\le C_{d,L}\sum_{k=0}^{L}t^{-k}|X_c|^k .
\]
Because \(L\ge d+1\) and \(|X_c|>t\),
\[
t^{-(d+1)}|X_c|^{d+1}\le t^{-L}|X_c|^L,
\qquad
t^{-k}|X_c|^k\le t^{-L}|X_c|^L\quad(0\le k\le L).
\]
Thus the first term of \(\mathcal R_{d,L}\) is the local Taylor remainder on
\(|X_c|\le t\), while its single tail term controls both the Poisson quotient
and every subtracted polynomial mode on \(|X_c|>t\).  The same pointwise
bound controls both \(L^\infty\) and \(L^1(\Omega_d)\), since \(\Omega_d\) is a
probability measure.

The assumption on \(\mathcal R_{d,L}\) is exactly the assertion that the
right side of \eqref{eq:multidim-density-quotient-explicit-remainder} is
\(o(t^{-L})\).
The constant term equals \(1\) and is removed by the definition of
\(\delta_t^{(d)}\).  The first-order Taylor layer is
\[
\sum_{|\alpha|=1}\mu_\alpha U_\alpha^{(d)}(y)t^{-1}=0
\]
because \(\nu\) is centred.  The remaining Taylor terms give
\eqref{eq:multidim-density-quotient-expansion-main}.

The uniform \(O(t^{-2})\) bound is a consequence of this same \(L\)-jet, not
of a separate order-two large-translation estimate.  Indeed, the coefficient
bound \eqref{eq:multidim-coefficient-bound} gives
\[
\sup_{y\in\RR^d}
\left|\sum_{2\le|\alpha|\le L}
\mu_\alpha U_\alpha^{(d)}(y)t^{-|\alpha|}\right|
\le C_{d,L,\nu}\sum_{k=2}^{L}t^{-k}
=O(t^{-2}),
\]
and \(\|\rho_{L,t}^{(d)}\|_\infty=o(t^{-L})=O(t^{-2})\) because \(L\ge2\).
Hence
\(\|\delta_t^{(d)}\|_\infty=O(t^{-2})\).
For all sufficiently large \(t\), therefore, \(\|\delta_t^{(d)}\|_\infty\le
1/2\), so \(1+\delta_t^{(d)}\) is bounded above and below by positive
constants.  Also
\[
\int_{\RR^d}\delta_t^{(d)}(y)\,\Omega_d(\dd y)
=\int_{\RR^d}\{h_t(\bar{\gamma}+ty)-g_t(\bar{\gamma}+ty)\}t^d\,\dd y=0,
\]
because both densities have total mass one.  Since \(\Phi\) is bounded by a
constant multiple of \(s^2\) on \(|s|\le1/2\), the integral
\(\int\Phi(\delta_t^{(d)})\,\dd\Omega_d\) is finite for all large \(t\), and
the change of variables \(x=\bar{\gamma}+ty\) gives the displayed KL
representation in \((\mathrm{EX}_{d,L})\).
\end{proof}

\begin{theorem}[Finite-moment supply of the quotient jet]%
\label{thm:multidimensional-density-quotient}
Let \(d\ge1\), let \(L\in\mathbb N\) satisfy \(L\ge d+1\), and let \(\nu\) be
a probability measure on \(\RR^d\) with mean \(\bar{\gamma}\) and finite
centred \(L\)th absolute moment.  Then
\(\mathcal R_{d,L}(t)=o(t^{-L})\), and the conclusions of
Proposition~\ref{thm:multidimensional-conditional-tail-density-quotient}
hold.
In particular \(\nu\) satisfies \((\mathrm{MPJ}_{d,L})\) and
\((\mathrm{EX}_{d,L})\) for all sufficiently large \(t\).
\end{theorem}

\begin{proof}
Lemma~\ref{lem:two-region-remainder-equivalent-finite-moment} gives
\(\mathcal R_{d,L}(t)=o(t^{-L})\).  Applying
Proposition~\ref{thm:multidimensional-conditional-tail-density-quotient} gives
the stated quotient jet and entropy-existence package.
\end{proof}

\begin{proposition}[Multidimensional quotient-to-KL transfer]%
\label{thm:multidimensional-poisson-entropy-jets}
Let \(d\ge1\), let \(L\in\mathbb N\) satisfy \(L\ge2\), let \(\nu\) be a
probability law on \(\RR^d\), and assume that \((\nu,d,L)\) satisfies both
\((\mathrm{MPJ}_{d,L})\) and \((\mathrm{EX}_{d,L})\).  Thus
\[
\delta_t^{(d)}(y)
=\sum_{2\le\lvert\alpha\rvert\le L}
\mu_\alpha U_\alpha^{(d)}(y)t^{-\lvert\alpha\rvert}
+\rho_{L,t}^{(d)}(y),
\]
with
\[
\lVert\rho_{L,t}^{(d)}\rVert_\infty+
\lVert\rho_{L,t}^{(d)}\rVert_{L^1(\Omega_d)}=o(t^{-L}),
\qquad
\lVert\delta_t^{(d)}\rVert_\infty=O(t^{-2}).
\]
Then
\[
D_{\mathrm{KL}}\bigl(P_t^{(d)}*\nu\,\|\,P_t^{(d)}(\cdot-\bar{\gamma})\bigr)
=\sum_{m=4}^{L+2}\mathcal A_m^{(d)}(\nu)t^{-m}
+o(t^{-(L+2)}),
\]
where
\[
\mathcal A_m^{(d)}(\nu)=
\sum_{r=2}^{\lfloor m/2\rfloor}
\frac{(-1)^r}{r(r-1)}
\sum_{\substack{\alpha_1,\ldots,\alpha_r\in\NN_0^d\\
|\alpha_1|+\cdots+|\alpha_r|=m\\
2\le|\alpha_j|\le L\ \mathrm{for}\ 1\le j\le r}}
\mu_{\alpha_1}\cdots\mu_{\alpha_r}
\int_{\RR^d}\prod_{j=1}^r
U_{\alpha_j}^{(d)}(y)\,\Omega_d(\dd y).
\]
Moreover, all odd retained coefficients vanish:
\[
\mathcal A_m^{(d)}(\nu)=0
\qquad\text{for every odd }m,\quad 4\le m\le L+2.
\]
\end{proposition}

\begin{remark}
For \(d=1\), \(\Sigma_0=0\), \(\operatorname{tr}\Sigma=\mu_2\), and the
quadratic form evaluated below gives \(c_1^{\rm iso}=1/8\).  Therefore, once
the quotient-jet hypothesis of this proposition is in force,
\(\mathcal A_4^{(1)}=A_4=\mu_2^2/8\), matching the one-dimensional
Cayley--Chebyshev coefficient.
\end{remark}

\begin{proof}
Write
\[
D_j^{(d)}(y)=\sum_{|\alpha|=j}\mu_\alpha U_\alpha^{(d)}(y),\qquad
d_t(y)=\sum_{j=2}^{L}D_j^{(d)}(y)t^{-j}.
\]
\((\mathrm{MPJ}_{d,L})\) gives
\[
\delta_t^{(d)}=d_t+\rho_{L,t}^{(d)},\qquad
\|d_t\|_\infty=O(t^{-2}),\qquad
\|\rho_{L,t}^{(d)}\|_\infty+
\|\rho_{L,t}^{(d)}\|_{L^1(\Omega_d)}=o(t^{-L}),
\]
and also \(\|\delta_t^{(d)}\|_\infty=O(t^{-2})\).  Since
\[
P_t^{(d)}(ty)t^d\,\dd y=P_1^{(d)}(y)\dd y=\Omega_d(\dd y)
\]
and \((\mathrm{EX}_{d,L})\) gives the large-\(t\) entropy representation
\[
D_{\rm KL}\bigl(P_t^{(d)}*\nu\,\|\,P_t^{(d)}(\cdot-\bar{\gamma})\bigr)
=\int_{\RR^d}\Phi(\delta_t^{(d)}(y))\,\Omega_d(\dd y),
\]
where \(\Phi(s)=(1+s)\log(1+s)-s\).
By Proposition~\ref{prop:finite-quotient-jet-implies-kl-coefficient-jet},
applied with \((E,\lambda)=(\RR^d,\Omega_d)\), \(J=L\), and
\(b_j=D_j^{(d)}\), only pure powers of the finite quotient polynomial
\(d_t\) contribute through order \(L+2\).  Collecting the coefficient of
\(t^{-m}\) in
\[
\sum_{r=2}^{\lfloor m/2\rfloor}
\frac{(-1)^r}{r(r-1)}d_t(y)^r
\]
gives precisely the displayed formula for \(\mathcal A_m^{(d)}(\nu)\), and
all omitted layers are \(o(t^{-(L+2)})\).  If \(m\) is odd, then every
product integral in the displayed coefficient has odd total mode order
\(m\), so Lemma~\ref{lem:multidim-cayley-poisson-mode-parity} makes it zero.
Thus all retained odd coefficients vanish.
\end{proof}

\begin{proposition}[Explicit-remainder form of the multidimensional KL jet]%
\label{thm:multidim-conditional-tail-poisson-kl-asymptotic}
Let \(d\ge1\), let \(L\in\mathbb N\) satisfy \(L\ge d+1\), and let \(\nu\) be
a probability measure on \(\RR^d\) with mean \(\bar{\gamma}\).  Put
\(X_c=\gamma-\bar{\gamma}\).  Assume
\(\mathcal R_{d,L}(t)=o(t^{-L})\); equivalently,
\(\EE|X_c|^L<\infty\).  Set
\[
\begin{aligned}
h_t&=P_t^{(d)}*\nu,\\
g_t&=P_t^{(d)}(\cdot-\bar{\gamma}),
\end{aligned}
\]
and define
\[
\mu_\alpha=\EE[X_c^\alpha]\qquad(|\alpha|\le L),
\]
which are finite by
Lemma~\ref{lem:two-region-remainder-equivalent-finite-moment}.  Then
\[
D_{\mathrm{KL}}(h_t\|g_t)
=\sum_{m=4}^{L+2}\mathcal A_m^{(d)}(\nu)t^{-m}
+o(t^{-(L+2)}),
\]
where
\[
\mathcal A_m^{(d)}(\nu)=
\sum_{r=2}^{\lfloor m/2\rfloor}
\frac{(-1)^r}{r(r-1)}
\sum_{\substack{\alpha_1,\ldots,\alpha_r\in\NN_0^d\\
|\alpha_1|+\cdots+|\alpha_r|=m\\
2\le|\alpha_j|\le L\ \mathrm{for}\ 1\le j\le r}}
\mu_{\alpha_1}\cdots\mu_{\alpha_r}
\int_{\RR^d}\prod_{j=1}^r
U_{\alpha_j}^{(d)}(y)\,\Omega_d(\dd y).
\]
All odd retained coefficients vanish:
\[
\mathcal A_m^{(d)}(\nu)=0
\qquad\text{for every odd }m,\quad 4\le m\le L+2.
\]
\end{proposition}

\begin{proof}
Proposition~\ref{thm:multidimensional-conditional-tail-density-quotient} proves
that the remainder condition \(\mathcal R_{d,L}(t)=o(t^{-L})\) implies
\((\mathrm{MPJ}_{d,L})+(\mathrm{EX}_{d,L})\), with density layers
\[
D_j^{(d)}(y;\nu)=\sum_{|\alpha|=j}\mu_\alpha U_\alpha^{(d)}(y).
\]
Proposition~\ref{thm:multidimensional-poisson-entropy-jets} applies to
exactly that quotient jet and gives the displayed entropy expansion and
coefficient formula.  The odd-coefficient cancellation is the
Lemma~\ref{lem:multidim-cayley-poisson-mode-parity} consequence recorded in
that proposition.  By
Lemma~\ref{lem:two-region-remainder-equivalent-finite-moment}, this is the
same hypothesis as finite \(L\)-absolute moment in the present range.
\end{proof}

\begin{corollary}[Finite-moment multidimensional KL jet]%
\label{thm:multidim-finite-moment-poisson-kl-asymptotic}%
\label{cor:multidim-finite-moment-entropy-jet}
Let \(d\ge1\), let \(L\in\mathbb N\) satisfy \(L\ge d+1\), and let \(\nu\) be
a probability measure on \(\RR^d\) with mean \(\bar{\gamma}\) and finite centred
\(L\)th absolute moment.  Set
\[
\begin{aligned}
h_t&=P_t^{(d)}*\nu,\\
g_t&=P_t^{(d)}(\cdot-\bar{\gamma}),
\end{aligned}
\]
and define the centred variable and moments by
\[
\begin{aligned}
X_c&=\gamma-\bar{\gamma},\\
\mu_\alpha&=\EE[X_c^\alpha]\qquad(|\alpha|\le L).
\end{aligned}
\]
Then
\[
D_{\mathrm{KL}}(h_t\|g_t)
=\sum_{m=4}^{L+2}\mathcal A_m^{(d)}(\nu)t^{-m}
+o(t^{-(L+2)}),
\]
where
\[
\mathcal A_m^{(d)}(\nu)=
\sum_{r=2}^{\lfloor m/2\rfloor}
\frac{(-1)^r}{r(r-1)}
\sum_{\substack{\alpha_1,\ldots,\alpha_r\in\NN_0^d\\
|\alpha_1|+\cdots+|\alpha_r|=m\\
2\le|\alpha_j|\le L\ \mathrm{for}\ 1\le j\le r}}
\mu_{\alpha_1}\cdots\mu_{\alpha_r}
\int_{\RR^d}\prod_{j=1}^r
U_{\alpha_j}^{(d)}(y)\,\Omega_d(\dd y).
\]
All odd retained coefficients vanish:
\[
\mathcal A_m^{(d)}(\nu)=0
\qquad\text{for every odd }m,\quad 4\le m\le L+2.
\]
\end{corollary}

\begin{proof}
Theorem~\ref{thm:multidimensional-density-quotient} proves that the finite
centred \(L\)th absolute moment hypothesis gives
\((\mathrm{MPJ}_{d,L})+(\mathrm{EX}_{d,L})\), with the explicit
\(\mathcal R_{d,L}\) remainder bound.  Applying
Proposition~\ref{thm:multidimensional-poisson-entropy-jets} gives the
displayed entropy expansion, coefficient formula, and, through
Lemma~\ref{lem:multidim-cayley-poisson-mode-parity}, parity cancellation.
No covariance contraction is used in this proof.
\end{proof}

\begin{remark}[Two coefficient routes]
Corollary~\ref{thm:multidim-finite-moment-poisson-kl-asymptotic} proves
the finite-moment multidimensional coefficient expansion.  Proposition
\ref{thm:multidim-conditional-tail-poisson-kl-asymptotic} is the same
finite-moment assertion written with the explicit remainder
\(\mathcal R_{d,L}\), whose \(o(t^{-L})\) decay is equivalent to finite
\(L\)-moment by
Lemma~\ref{lem:two-region-remainder-equivalent-finite-moment}.  The leading
coefficient route is the one stated at the start of this subsection.
\end{remark}
